% from ega0-19-fr.tex
\oldpage[0\textsubscript{IV-1}]{95}

\begin{env}[19.5.6.2]
\label{0.19.5.6.2-fr}
Pour tout entier $n\geq0$, nous poserons
\[
  E_n=B/\mathfrak J^{n+1},
  \qquad\text{de sorte que}\qquad E_0=C~;
\]
les idéaux $(\mathfrak b_\alpha+\mathfrak J^{n+1})/\mathfrak J^{n+1}$ forment
un système fondamental d'idéaux ouverts dans $E_n$, et nous poserons
\[
  E_{\alpha,n}=B/(\mathfrak b_\alpha+\mathfrak J^{n+1})
  =B_\alpha/\mathfrak J_\alpha^{n+1},
\]
anneau quotient discret. De même, dans
$\gr_{\mathfrak J}^n(B)=\mathfrak J^n/\mathfrak J^{n+1}$, nous avons vu que les
$((\mathfrak b_\alpha\cap\mathfrak J^n)+\mathfrak J^{n+1})/
\mathfrak J^{n+1}$ forment un système fondamental de voisinages de $0$, le
quotient de $\mathfrak J^n/\mathfrak J^{n+1}$ par ce sous-module s'identifiant
canoniquement à $\mathfrak J_\alpha^n/\mathfrak J_\alpha^{n+1}$. Considérons
l'algèbre symétrique
$\mathbf S_{C_\alpha}^*(\mathfrak J_\alpha/\mathfrak J_\alpha^2)$~; nous
noterons $F_{\alpha,n}$ le quotient de
$\mathbf S_{C_\alpha}^*(\mathfrak J_\alpha/\mathfrak J_\alpha^2)$ par la
puissance $(n+1)$-ième de son idéal d'augmentation. Pour un $n\geq1$ fixé, il
résulte de (19.5.1.1) que les
$\mathbf S_{C_\alpha}^n(\mathfrak J_\alpha/\mathfrak J_\alpha^2)$ sont les
quotients de $\mathbf S_C^n(\mathfrak J/\mathfrak J^2)$ par un système
fondamental de sous-modules ouverts dans ce $C$-module topologique.

Pour abréger le langage, nous dirons que pour $\alpha\leq\beta$, les
homomorphismes canoniques $B_\beta\to B_\alpha$, $C_\beta\to C_\alpha$,
$\mathfrak J_\beta/\mathfrak J_\beta^2\to
\mathfrak J_\alpha/\mathfrak J_\alpha^2$, $E_{\beta,n}\to E_{\alpha,n}$,
$F_{\beta,n}\to F_{\alpha,n}$, etc., sont les \emph{homomorphismes de
transition}.
\end{env}

\begin{lemma}[19.5.6.3]
\label{0.19.5.6.3-fr}
Supposons que la $A$-algèbre $C$ soit formellement lisse, l'anneau $B$
préadmissible et le $C$-module $\mathfrak J/\mathfrak J^2$ formellement
projectif. Alors~:
\begin{enumerate}
  \item[\rm(i)] Pour tout $\alpha$, il existe $\beta\geq\alpha$ et un
  $A$-homomorphisme surjectif d'algèbres
  \[
    v_{\alpha\beta}:F_{\beta,n}\longrightarrow E_{\alpha,n}
  \]
  tel que $\gr^0(v_{\alpha\beta}):C_\beta\to C_\alpha$ et
  $\gr^1(v_{\alpha\beta}):\mathfrak J_\beta/\mathfrak J_\beta^2
  \to\mathfrak J_\alpha/\mathfrak J_\alpha^2$ soient les homomorphismes de
  transition.
  \item[\rm(ii)] Si $\beta$ et $v_{\alpha\beta}$ vérifient les conditions de
  \emph{(i)}, alors, pour tout $\gamma\geq\beta$, il existe
  $\delta\geq\gamma$ et un $A$-homomorphisme surjectif d'algèbres
  $v_{\gamma\delta}:F_{\delta,n}\to E_{\gamma,n}$ vérifiant les conditions de
  \emph{(i)} (pour $\gamma$ et $\delta$) et rendant commutatif le diagramme
  \begin{equation}
  \label{0.19.5.6.4-fr}
  \xymatrix@C=3em{
    F_{\beta,n} \ar[r]^{v_{\alpha\beta}} & E_{\alpha,n} \\
    F_{\delta,n} \ar[u] \ar[r]_{v_{\gamma\delta}} & E_{\gamma,n} \ar[u]
  }
  \tag{19.5.6.4}
  \end{equation}
  où les flèches verticales sont les homomorphismes de transition.
\end{enumerate}
\end{lemma}

\begin{env}[19.5.6.4]
\label{0.19.5.6.4-proof-fr}
\begin{enumerate}
  \item[\rm(i)] Dans la $A$-algèbre topologique discrète $E_{\alpha,n}$,
  l'idéal $\mathfrak J_\alpha/\mathfrak J_\alpha^{n+1}$ est nilpotent, et
  l'isomorphisme identique
  $C_\alpha\xrightarrow{\sim}
  E_{\alpha,n}/(\mathfrak J_\alpha/\mathfrak J_\alpha^{n+1})$ donne par
  composition un $A$-homomorphisme continu
  $C\to C_\alpha\to
  E_{\alpha,n}/(\mathfrak J_\alpha/\mathfrak J_\alpha^{n+1})$~;

